% CXL内存扩展架构图
% 展示Type 1/2/3设备类型和CXL拓扑

\documentclass[tikz,border=10pt]{standalone}
\usepackage{tikz}
\usepackage{ctex}
\usetikzlibrary{shapes.geometric, arrows.meta, positioning, calc, fit, backgrounds, decorations.pathreplacing}

\begin{document}
\begin{tikzpicture}[
    node distance=0.4cm,
    box/.style={rectangle, rounded corners, draw, minimum width=2.5cm, minimum height=0.8cm, align=center, font=\scriptsize},
    cpu/.style={box, fill=blue!20},
    device/.style={box, fill=green!20},
    switch/.style={box, fill=yellow!20, minimum width=1.5cm},
    memory/.style={cylinder, draw, shape border rotate=90, aspect=0.3, fill=orange!20, minimum width=1.5cm, minimum height=0.8cm, align=center, font=\scriptsize},
    arrow/.style={-{Stealth[length=2mm]}, thick},
    label/.style={font=\scriptsize\bfseries}
]

% 标题
\node[font=\large\bfseries] at (8, 7) {CXL内存扩展架构};

% ============ 主机CPU ============
\node[cpu, minimum width=3cm, minimum height=1.2cm] (cpu) at (4, 5) {\textbf{主机CPU}\\CXL Root Complex};

% CXL协议栈
\node[below=0.3cm of cpu, font=\tiny, text width=3cm, align=center, draw, fill=blue!10, rounded corners] {
    CXL.io | CXL.cache | CXL.mem
};

% ============ Type 1 设备 (网卡) ============
\node[device] (type1) at (0, 2.5) {\textbf{Type 1}\\网卡/NIC};
\node[below=0.1cm of type1, font=\tiny, text width=2.5cm, align=center] {
    缓存一致性设备\\CXL.cache + CXL.io
};

% ============ Type 2 设备 (GPU) ============
\node[device] (type2) at (4, 2.5) {\textbf{Type 2}\\GPU/加速器};
\node[memory] (hbm) at (4, 1) {HBM\\显存};
\node[below=0.1cm of hbm, font=\tiny, text width=2.5cm, align=center] {
    带内存加速器\\CXL.cache + CXL.mem
};

% ============ Type 3 设备 (内存扩展) ============
\node[device] (type3) at (8, 2.5) {\textbf{Type 3}\\内存扩展卡};
\node[memory] (ddr) at (8, 1) {DDR5\\内存};
\node[below=0.1cm of ddr, font=\tiny, text width=2.5cm, align=center] {
    纯内存设备\\CXL.mem + CXL.io
};

% ============ CXL交换机与内存池化 ============
\node[switch, minimum width=2cm, minimum height=1cm] (cxl-sw) at (12, 4) {\textbf{CXL}\\\textbf{交换机}};

% 内存池
\node[memory] (pool1) at (11, 2) {内存池 A};
\node[memory] (pool2) at (13, 2) {内存池 B};
\node[below=0.1cm of pool1, font=\tiny] {共享内存};
\node[below=0.1cm of pool2, font=\tiny] {共享内存};

% ============ 连接线路 ============
% CPU到设备
\draw[arrow, blue!50] (cpu) -- (type1) node[midway, left, font=\tiny] {CXL.io};
\draw[arrow, blue!50] (cpu) -- (type2) node[midway, right, font=\tiny] {CXL.cache};
\draw[arrow, blue!50] (cpu) -- (type3) node[midway, right, font=\tiny] {CXL.mem};

% GPU到HBM
\draw[arrow, orange!50] (type2) -- (hbm);

% 内存扩展卡到DDR
\draw[arrow, orange!50] (type3) -- (ddr);

% CPU到交换机
\draw[arrow, green!50] (cpu) -- (cxl-sw) node[midway, above, font=\tiny] {CXL Link};

% 交换机到内存池
\draw[arrow, green!50] (cxl-sw) -- (pool1);
\draw[arrow, green!50] (cxl-sw) -- (pool2);

% ============ 右侧：CXL 3.0 Fabric拓扑 ============
\node[font=\bfseries] at (16, 5.5) {CXL 3.0 Fabric拓扑};

% 多级交换
\node[switch] (sw1) at (15, 4.5) {交换机 L1};
\node[switch] (sw2) at (17, 4.5) {交换机 L1};
\node[switch] (sw-root) at (16, 3.5) {根交换机};

\node[cpu, minimum width=1.5cm] (host1) at (14.5, 2.5) {主机1};
\node[cpu, minimum width=1.5cm] (host2) at (15.5, 2.5) {主机2};
\node[cpu, minimum width=1.5cm] (host3) at (16.5, 2.5) {主机3};
\node[cpu, minimum width=1.5cm] (host4) at (17.5, 2.5) {主机4};

\node[memory] (mem1) at (14.5, 1.5) {内存};
\node[memory] (mem2) at (17.5, 1.5) {内存};

% 连接
\draw[arrow, gray] (sw1) -- (sw-root);
\draw[arrow, gray] (sw2) -- (sw-root);
\draw[arrow, gray] (host1) -- (sw1);
\draw[arrow, gray] (host2) -- (sw1);
\draw[arrow, gray] (host3) -- (sw2);
\draw[arrow, gray] (host4) -- (sw2);
\draw[arrow, gray] (mem1) -- (host1);
\draw[arrow, gray] (mem2) -- (host4);

\node[below=0.5cm of host3, font=\tiny, text width=4cm, align=center] {
    支持最多4096个节点\\全局内存共享
};

% ============ 协议栈说明 ============
\node[below=4.5cm of cpu, font=\scriptsize, text width=14cm, align=left, draw, rounded corners, fill=gray!10] {
    \textbf{CXL协议栈：}\\[3pt]
    \begin{tabular}{ll}
    \textcolor{blue!60!black}{CXL.io} & 兼容PCIe协议，用于设备配置、寄存器访问和传统I/O操作 \\
    \textcolor{green!60!black}{CXL.cache} & 支持设备缓存主机内存，维护缓存一致性（类似CPU缓存一致性协议） \\
    \textcolor{orange!60!black}{CXL.mem} & 主机直接访问设备内存，使用普通load/store指令 \\
    \end{tabular}
};

% ============ 应用场景 ============
\node[below=6cm of cpu, font=\scriptsize, text width=14cm, align=left, draw, rounded corners, fill=blue!5] {
    \textbf{典型应用场景：}\\[3pt]
    \begin{tabular}{lll}
    \textbf{Type 1} & 智能网卡 & 网卡缓存CPU内存数据，实现零拷贝网络 \\
    \textbf{Type 2} & AI加速器 & GPU与CPU共享内存空间，消除数据拷贝开销 \\
    \textbf{Type 3} & 内存扩展 & 突破主板DIMM插槽限制，支持TB级内存扩展 \\
    \textbf{Fabric} & 内存池化 & 多台服务器共享内存池，动态分配，提高利用率 \\
    \end{tabular}
};

% ============ 版本演进 ============
\node[right=0.5cm of cxl-sw, font=\tiny, text width=3cm, align=left, draw, fill=yellow!10, rounded corners] {
    \textbf{CXL演进：}\\
    1.0/1.1: 基础互联\\
    2.0: 内存池化\\
    3.0: Fabric能力\\
    3.1: 点对点通信
};

\end{tikzpicture}
\end{document}
